\chapter{The Female Reproductive System}
\index{The Female Reproductive System}

\section*{Definition and Overview}
The female reproductive system is the group of organs that produces eggs (ova), enables fertilisation and pregnancy, and supports a developing baby until birth. It includes the ovaries, fallopian tubes, uterus, cervix, vagina, and vulva, along with the hormonal system that regulates their function. The system also produces the female sex hormones oestrogen and progesterone, which shape sexual development, the menstrual cycle, and many other aspects of health. Understanding how these organs work is the foundation for understanding menstruation, fertility, contraception, pregnancy, and the many conditions that can affect them. This chapter provides an overview of the anatomy and function of the female reproductive system.

\section*{Epidemiology}
Disorders of the female reproductive system are common across the lifespan. Around one in four women experiences heavy menstrual bleeding, up to one in nine has endometriosis, and polycystic ovary syndrome affects 8--13\% of women of reproductive age. Infertility affects roughly one in six couples. Gynaecological cancers, including cervical, endometrial, and ovarian cancer, affect thousands of women each year, though cervical cancer has fallen dramatically with screening and HPV vaccination. Menstrual and reproductive health problems are among the most common reasons for GP visits among young women, and their impact on work, study, and quality of life is substantial.

\section*{Aetiology and Pathophysiology}
The system is organised into organs with distinct roles in reproduction and hormonal regulation.
\begin{itemize}
    \item \textbf{Ovaries:} produce eggs within follicles and secrete oestrogen, progesterone, and small amounts of testosterone; a woman is born with all the eggs she will ever have, around one to two million, which decline to a few hundred thousand by puberty
    \item \textbf{Fallopian tubes:} transport the egg towards the uterus; fertilisation normally occurs here
    \item \textbf{Uterus (womb):} a muscular organ whose lining (endometrium) thickens each cycle to receive a fertilised egg, and which expands greatly in pregnancy; its muscle powers labour
    \item \textbf{Cervix:} the lower neck of the uterus, which produces mucus that changes across the cycle and dilates in labour
    \item \textbf{Vagina:} the muscular canal from the cervix to the vulva; the birth canal and the organ of intercourse
    \item \textbf{Vulva:} the external genitals, including the labia and clitoris
    \item \textbf{Hormonal control:} the hypothalamus and pituitary gland direct the ovaries through follicle-stimulating hormone (FSH) and luteinising hormone (LH), driving the menstrual cycle
    \item \textbf{The menstrual cycle:} approximately monthly, with a follicular phase of egg maturation, ovulation, and a luteal phase preparing the endometrium; if no pregnancy occurs, the lining is shed as a period
\end{itemize}

\section*{Clinical Features}
Normal function and common disorders produce recognisable patterns.
\begin{itemize}
    \item Menstruation: regular monthly bleeding of about 2--7 days, normally pain-free or mildly uncomfortable
    \item Ovulation: mid-cycle release of an egg, with occasional one-sided twinges (mittelschmerz)
    \item Fertility: the ability to conceive, with most couples conceiving within 12 months of trying
    \item Pregnancy: the uterus supporting a fetus for about 40 weeks, ending in labour and delivery
    \item Menopause: the permanent end of menstruation around age 51, marking the end of reproductive capacity
    \item Common symptoms of disorders: heavy, painful, irregular, or absent periods; pelvic pain; abnormal discharge; and infertility
    \item Hormonal effects: changes in libido, mood, skin, and bone across the cycle and the lifespan
\end{itemize}

\section*{Differential Diagnosis}
Symptoms such as pelvic pain and abnormal bleeding have many possible sources.
\begin{itemize}
    \item Structural conditions: fibroids, polyps, endometriosis, adenomyosis, and ovarian cysts
    \item Hormonal disorders: polycystic ovary syndrome, thyroid disease, and premature ovarian insufficiency
    \item Infections: sexually transmitted infections and pelvic inflammatory disease
    \item Pregnancy-related conditions: ectopic pregnancy and miscarriage
    \item Cancer: cervical, endometrial, and ovarian cancer, especially in older women
    \item Functional disorders: primary dysmenorrhoea and dysfunctional uterine bleeding
    \item Urinary and bowel conditions that cause pelvic symptoms
\end{itemize}

\section*{When to Seek Medical Care}
Women should seek care for new or unusual bleeding, severe period pain, persistent pelvic pain, unusual discharge, or concerns about fertility. Regular cervical screening and preventive visits are part of maintaining reproductive health.

\textbf{Contact your local emergency services for:}
\begin{itemize}
    \item Very heavy vaginal bleeding that soaks pads within an hour, especially with faintness
    \item Sudden severe abdominal or pelvic pain
    \item Bleeding with pain in early pregnancy, which may indicate ectopic pregnancy
    \item Fever with pelvic pain, which may indicate pelvic infection or sepsis
    \item Collapse or loss of consciousness
\end{itemize}

\section*{Diagnosis and Investigations}
\begin{itemize}
    \item \textbf{History:} menstrual, obstetric, gynaecological, sexual, and contraceptive history
    \item \textbf{Examination:} abdominal and pelvic examination (speculum and bimanual) as appropriate
    \item \textbf{Cervical screening:} HPV test every 5 years from age 25 to 74
    \item \textbf{Ultrasound:} transvaginal or transabdominal imaging of the uterus, endometrium, and ovaries
    \item \textbf{Blood tests:} pregnancy test, hormone profiles, full blood count, and iron studies
    \item \textbf{Endometrial sampling and hysteroscopy:} for abnormal bleeding, particularly after menopause
    \item \textbf{Laparoscopy:} for suspected endometriosis and other pelvic conditions
\end{itemize}

\section*{Management}
\begin{itemize}
    \item \textbf{Menstrual disorders:} hormonal treatments, pain relief, intrauterine devices, and surgery including endometrial ablation and hysterectomy for selected cases
    \item \textbf{Endometriosis:} hormonal suppression, analgesia, and laparoscopic surgery, managed by a multidisciplinary team
    \item \textbf{Infections:} antibiotics for STIs and pelvic inflammatory disease, with partner treatment
    \item \textbf{Contraception and family planning:} a full range of hormonal, barrier, and permanent methods
    \item \textbf{Fertility:} investigation and treatment including ovulation induction and IVF
    \item \textbf{Gynaecological cancer:} surgery, radiotherapy, and chemotherapy through specialist services
    \item \textbf{Menopause:} lifestyle measures, hormone therapy after discussion, and vaginal treatments
    \item \textbf{Preventive care:} HPV vaccination, cervical screening, and breast screening
\end{itemize}

\section*{Complications}
\begin{itemize}
    \item Chronic pelvic pain and reduced quality of life from endometriosis and other conditions
    \item Anaemia from heavy menstrual bleeding
    \item Infertility from tubal damage, endometriosis, and hormonal disorders
    \item Ectopic pregnancy rupture, a life-threatening emergency
    \item Pelvic infection progressing to sepsis
    \item Gynaecological cancers, with prognosis improved by early detection
    \item Pelvic organ prolapse and incontinence in later life
\end{itemize}

\section*{Prognosis and Outcomes}
Most disorders of the female reproductive system are manageable, and many are curable. Menstrual disorders respond well to medical and surgical treatment; endometriosis is a chronic condition that can be effectively controlled; and most couples with infertility conceive with treatment. Screening has transformed the outcomes of cervical cancer, and survival from gynaecological cancers continues to improve. With access to good preventive care, the female reproductive system functions normally through its lifespan, and its disorders are recognised and treated early.

\section*{Prevention and Self-Care}
\begin{itemize}
    \item Attend cervical screening every 5 years from age 25
    \item Have the HPV vaccine, ideally before sexual debut
    \item Use condoms to prevent sexually transmitted infections
    \item Maintain a healthy weight, exercise, and eat well, which supports hormonal health
    \item Track periods to notice changes early
    \item See a doctor for pain, bleeding, or other symptoms rather than enduring them
    \item Plan contraception and pregnancy with medical advice
    \item Seek pelvic floor physiotherapy for prolapse or incontinence
\end{itemize}

\section*{Sources and Further Reading}
The following reputable sources provide further information for healthcare students and patients seeking reliable, current advice about the female reproductive system.
\begin{itemize}
    \item Healthdirect Australia. \textit{Female reproductive system.} https://www.healthdirect.gov.au/
    \item Better Health Channel. \textit{The female reproductive system.} https://www.betterhealth.vic.gov.au/
    \item Jean Hailes for Women's Health. \textit{Reproductive health.} https://www.jeanhailes.org.au/
    \item NHS. \textit{Female reproductive system.} https://www.nhs.uk/
    \item Mayo Clinic. \textit{Female reproductive system.} https://www.mayoclinic.org/
    \item MedlinePlus. \textit{Female reproductive system.} https://medlineplus.gov/
\end{itemize}
